% exonot-demo.tex -- demonstration of the exonot package.
% Compile with:  pdflatex exonot-demo.tex
% Dependencies (all standard TeX Live): amsmath, amssymb, booktabs, geometry.

\documentclass{article}

\usepackage{amsmath}
\usepackage{amssymb}
\usepackage{booktabs}
\usepackage[margin=1in]{geometry}

\usepackage[symbols]{exonot}

\title{The \texttt{exonot} package: standardized exoplanet notation}
\author{Rishabh Garg}
\date{2026-06-27}

\begin{document}
\maketitle

\begin{abstract}
We present transmission spectroscopy of the hot Saturn
\system{WASP-39}{b}, a \ensuremath{0.28}\,\Mjup{} planet orbiting its
host \star{A} with a period \Porb${}=4.06$\,d at a semi-major axis
\aorb${}=0.0486$\,au.  The retrieved spectrum shows strong \HO{},
\COtwo, and \SOtwo{} absorption, with tentative \CO{} and \CH, in a
\Htwo/\Helium{} dominated atmosphere of mean molecular weight
\mmw${}\approx2.3$ and carbon-to-oxygen ratio \CtoO${}\approx0.35$.  The
inferred equilibrium temperature is \Teq${}\approx1170$\,K and the
atmospheric scale height \Hatm{} is large enough to yield a transit
signal of several hundred \ppm{}, detected at \snr${}>20$ per
resolution element.
\end{abstract}

\section{System and stellar parameters}

The host star \star{A} has effective temperature \Teff, surface gravity
\logg, luminosity \Lsun, mass \Mstar${}=0.93$\,\Msun, radius
\Rstar${}=0.90$\,\Rsun, and metallicity \FeH${}=-0.10$.  The orbit is
consistent with circular (\ecc${}<0.05$) and well aligned
(\obliq${}\approx0^\circ$), with inclination \inc${}=87.8^\circ$ and
radial-velocity semi-amplitude \Krv${}=33$\,m\,s$^{-1}$.  Planetary
parameters for \system{WASP-39}{b} are collected in
Table~\ref{tab:params}.  Companion candidates would be labelled
\planet{c}, \planet{d}, and so on.

\begin{table}[ht]
  \centering
  \caption{Measured and derived parameters of \system{WASP-39}{b}.}
  \label{tab:params}
  \begin{tabular}{llr}
    \toprule
    Parameter                 & Symbol     & Value            \\
    \midrule
    Planet mass               & \Mp        & $0.28$\,\Mjup    \\
    Planet radius             & \Rp        & $1.27$\,\Rjup    \\
    Planet density            & \rhop      & $0.17$\,g\,cm$^{-3}$ \\
    Orbital period            & \Porb      & $4.06$\,d        \\
    Semi-major axis           & \aorb      & $0.0486$\,au     \\
    Eccentricity              & \ecc       & $<0.05$          \\
    RV semi-amplitude         & \Krv       & $33$\,m\,s$^{-1}$\\
    Equilibrium temperature   & \Teq       & $1170$\,K        \\
    Surface gravity           & \logg      & $2.63$           \\
    Impact parameter          & \bparam    & $0.45$           \\
    Transit depth             & \tdepth    & $0.0211$         \\
    Total transit duration    & \tdur      & $2.80$\,h        \\
    Ingress/egress duration   & \ting      & $0.24$\,h        \\
    Limb-darkening coeff.     & \ldc[1]    & $0.42$           \\
    Mean molecular weight     & \mmw       & $2.3$            \\
    \bottomrule
  \end{tabular}
\end{table}

\section{Atmospheric scale height}

The pressure scale height that sets the amplitude of the transmission
signal follows from hydrostatic equilibrium,
\begin{equation}
  \Hatm = \frac{\kB\,\Teq}{\mmw\,\mH\,g},
  \label{eq:scaleheight}
\end{equation}
where \kB{} is the Boltzmann constant, \mH{} the mass of a hydrogen
atom, and $g$ the surface gravity (here written as \logg).  The
corresponding transit depth is
\begin{equation}
  \delta = \tdepth,
\end{equation}
so the spectral modulation across the transmission spectrum \Rspec{}
scales as $2\Hatm/R_{\starsym}$ in units of \tdepth.

\section{Composition}

Retrieved volume mixing ratios for the dominant absorbers are
\vmr{\HO}${}\sim10^{-3}$, \vmr{\COtwo}${}\sim10^{-4}$,
\vmr{\CO}${}\sim10^{-4}$, and \vmr{\CH}${}<10^{-6}$, alongside a
photochemically produced \SOtwo{} component and high-altitude hazes.
Optical absorbers such as \NaI, \KI, \TiO, and \VO{} are not required by
the data, nor are \Htwo{}O-dissociation tracers like \Hminus{} or
\eminus.  Model comparison favours the cloudy, \SOtwo-bearing fit with
$\Delta\BIC\approx-32$ (equivalently a Bayes factor $\ln\Bfac\approx16$
and \redchisq${}\approx1.05$).  After correcting for a stellar
contamination factor \fcont${}\approx0.99$, the data remain consistent
with a \Htwo/\Helium{} envelope of metallicity \FeH${}\approx+1$.  Earth
(\Earth) and Sun (\Sun) symbols are available when the \texttt{symbols}
option is loaded.

\end{document}
